\documentclass[a4paper,11pt]{report}
%\usepackage[utf8]{inputenc}



%\usepackage[justification=centerlast,font=small]{caption}
\usepackage[T1]{fontenc}

% amssymb musi być przed \usepackage[polish]{babel}
% żeby nie było błędu ``command \lll already defined''
\usepackage{amssymb} 
                     
\usepackage[polish]{babel}
\usepackage{polski}

%\usepackage[utf8]{inputenc}
\usepackage{wrapfig}
%\usepackage{dirtytalk}
%\usepackage[demo]{graphicx}
\usepackage{siunitx}
\usepackage{textcomp}
\usepackage{icomma}
\usepackage{threeparttable}
\usepackage{graphicx}
\usepackage{subfig}
\usepackage{setspace}
%\usepackage{pdflscape}
\usepackage[paper=portrait,pagesize]{typearea} % dla stron obróconych
%% dla malych arrow
\usepackage{stmaryrd}
\newcolumntype{C}{>{\centering\arraybackslash}X}
\usepackage{longtable}
\usepackage{gensymb}
%\usepackage{hyperref}

% podpis tabeli nad tabelą
\usepackage{floatrow}
\floatstyle{plaintop}
\restylefloat{table}

\usepackage[bottom]{footmisc}
\usepackage[figurename=Fig.]{caption}
%\usepackage[section]{placeins}
\usepackage{color, soul} % highlighting
\newcolumntype{Y}{>{\centering\arraybackslash}X} %necessary for tabularx to make equal width columns
\usepackage[usenames,dvipsnames]{xcolor} 
\usepackage[colorinlistoftodos,prependcaption,textsize=tiny]{todonotes}
%\newcommand{\hlc}[2][yellow]{ {\sethlcolor{#1} \hl{#2}} }
%\usepackage{graphicx} % for improved inclusion of graphics
%\usepackage{epstopdf}
%\usepackage{indentfirst} % wstawia wciecie w pierwszym akapicie rozdzia³u 
%\usepackage{amsfonts} % for \mathbf
%\usepackage{datetime}%zeby miec godzine kompilacji (opcja \date i \currenttime)
\usepackage{tabularx}%do robienia tabel
\usepackage{float}
\usepackage{amsmath}
%\usepackage{subcaption}
\usepackage{multirow}
\usepackage{booktabs}
\selectlanguage{polish}
%%%%%%%%%%%%%%%%%%%%%%% to prevent hyphenation %%%%%%%%%%%%%%%%%%%%%%%
%\usepackage[polish=nohyphenation]{hyphsubst}
\tolerance=1
\emergencystretch=\maxdimen
\hyphenpenalty=10000
\hbadness=10000
%\hyphenation{sze-regowo}
%%%%%%%%%%%%%%%%%%%%%%% end of to prevent hyphenation %%%%%%%%%%%%%%%%%%%%%%%
\usepackage{caption}% zamienia Rysunek na Rys. w podpisach rysunków
\captionsetup{figurename=Rys.}
\captionsetup{tablename=Tab.}
\captionsetup[table]{labelsep=space} % removes colon from caption
\captionsetup[figure]{labelsep=space} % removes colon from caption
% Zdefiniowanie autora i~tytułu:
\frenchspacing %wyłączenie wstawiania wiêkszych odstêpów na koñcu zdañ
\addtolength{\hoffset}{-0.5cm}
\addtolength{\textwidth}{1.0cm}%zwieksza szer. szpalty o liczbe cm
\usepackage{fancyhdr} %wstawia naglowki na stronach
%\pagestyle{headings}

%\usepackage[draft]{todonotes}   % notes showed; should be after xcolor

% do kolorowania pojedynczych komórek w tabelach \cellcolor
\usepackage{xcolor,colortbl}

\setcounter{tocdepth}{3} %%%% wyrzuca subsubsections z toc
\usepackage[most]{tcolorbox} % do rysowania kolorowych ramek z tekstem
\usepackage{hyperref}%do hiperłączy; powinno być jako ostatnie w preambule
%\hypersetup{colorlinks=false}%sprawia że kolorowe ramki hiperłączy nie są drukowane
\hypersetup{colorlinks,citecolor=black,filecolor=black,linkcolor=black,urlcolor=red}%wyłącza kolorowe ramki hiperłączy; linkiw www sa na czerwono

\usepackage{listings}
\usepackage{xcolor}
\usepackage{indentfirst} % wciecie akapitu

%New colors defined below
\definecolor{codegreen}{rgb}{0,0.6,0}
\definecolor{codegray}{rgb}{0.5,0.5,0.5}
\definecolor{codepurple}{rgb}{0.58,0,0.82}
\definecolor{backcolour}{rgb}{0.95,0.95,0.92}

%Code listing style named "mystyle"
\lstdefinestyle{mystyle}{
  backgroundcolor=\color{backcolour},   commentstyle=\color{codegreen},
  keywordstyle=\color{magenta},
  numberstyle=\tiny\color{codegray},
  stringstyle=\color{codepurple},
  basicstyle=\ttfamily\footnotesize,
  breakatwhitespace=false,         
  breaklines=true,                 
  captionpos=b,                    
  keepspaces=true,                 
  numbers=left,                    
  numbersep=5pt,                  
  showspaces=false,                
  showstringspaces=false,
  showtabs=false,                  
  tabsize=2
}

%"mystyle" code listing set
\lstset{style=mystyle}


\begin{document}
%\pagenumbering{Roman}
% Wstawienie autora i~tytu³u do sk³adu:
\begin{titlepage}
\begin{flushright}
\underline{Na prawach rękopisu}
\end{flushright}
 \vspace{0.5cm}
\begin{center}

\includegraphics[width=0.5\textwidth]{figures/logo.png}\\
{\large{POLITECHNIKA WROCŁAWSKA}}\\
{\large{WYDZIAŁ MECHANICZNO-ENERGETYCZNY}}
 \vspace{5cm}

 {\LARGE\bfseries ,,Szablon''\\}
 % ----------------------------------------------------------------
 \vspace{1.5cm}
 {\Large \textsc{dr inż. Przemysław Błasiak}}\\[5pt]
 
  % ----------------------------------------------------------------
 \vspace{1.0cm}
% \begin{table}[H]
%\begin{tabularx}{\linewidth}{YYYY}\\
%
%\multicolumn{1}{c}{}   & \multicolumn{1}{c}{}    & \multicolumn{1}{r}{Słowa kluczowe:} & \multicolumn{1}{l}{intensyfikacja wymiany ciepła}  \\
%\multicolumn{1}{c}{}   & \multicolumn{1}{c}{}    &                 & \multicolumn{1}{l}{pulsacyjna rurka ciepła}  \\
%\multicolumn{1}{c}{}   & \multicolumn{1}{c}{}    &                 & \multicolumn{1}{l}{przepływ wielofazowy}  \\
%
%\end{tabularx}
%\end{table}
%
%\vspace{0.5cm}
%\textsc{\large{{\textbf{Praca doktorska}}}} \\[5pt]
{\Large{Start pisania: 06.07.2021 r.}}\\[5pt]
\vspace{6.0cm}
% {By}\\[5pt] {\Large \sc {Me}}
 %\vfill
 % ----------------------------------------------------------------
%\includegraphics[width=0.19\textwidth]{example-image-a}\\[5pt]
%{\Large{Promotor pomocniczy: dr inż. Przemysław Błasiak}}\\[5pt]

%\vspace{0.5cm}
{Wrocław, 2021}
\end{center}
\end{titlepage}

\tableofcontents




\addcontentsline{toc}{chapter}{Bibliografia}
\bibliographystyle{plain}
\bibliography{szablonRefs}
\lstlistoflistings

\end{document}
